\chapter*{Executive Summary}
\addcontentsline{toc}{chapter}{Executive Summary}
\textbf{Problem and motivation.} Retail stores record CCTV footage mainly for security, yet managers lack a low-cost way to know whether customers are attended to, how staff time is distributed, and which zones receive traffic. Commercial retail-analytics products typically depend on dedicated sensors or hardware and on subscription costs, which puts them out of reach of many small and mid-sized retailers.

\textbf{Proposed solution.} RetailPulse reuses existing CCTV infrastructure. Vision components for Detection, Classification, Engagement and Trajectory convert video into per-entity analytics: footfall, approach time, interaction duration, idle time, unattended-customer flags, dwell time, paths, and heatmaps. Results are presented on a web dashboard for managers and a focused workspace for employees.

\textbf{Scope.} The committed scope is a single-store prototype with three cameras, running on CPU only. Multi-store scaling, full cross-camera re-identification, and shelf-level interaction detection are deliberately deferred to future work after a feasibility triage against hardware, timeline, and dependency constraints.

\textbf{Current status.} The Next.js frontend covers the manager, employee, and shared pages. FastAPI and SQLite support account, camera, alert and operational-task workflows. Role-aware access is enforced with signed HttpOnly session cookies and an internal secret header. A local interaction pipeline combines a trained employee/customer detector, BoT-SORT tracking, body-height proximity with a pass-by filter, and a YOLO pose judge that confirms an interaction only when someone faces the other person. A metric floor-distance prototype was built and then descoped because the moving-camera footage supported almost no floor projections. The dashboard vision metrics still partly use prototype or demo data.

\textbf{Validation, risks, and remaining work.} The detector reaches mAP50 0.969 on 455 held-out images. On the one labelled clip the pose judge raises interaction precision from 0.50 to 1.00, and the pipeline completes all 11 local clips; 44 interaction unit tests pass. Accuracy is measured on one labelled clip only. The main risks are CPU-only inference speed, role-classification accuracy, depth ambiguity in a single image, and identity switches on moving footage. The remaining work covers vision-agent integration, evaluation on store footage, hardening, and the final demonstration.
